\documentclass[11pt]{article}
\usepackage[margin=1in]{geometry}
\usepackage{amsmath,amssymb,amsthm}
\usepackage{xurl}
\usepackage{hyperref}
\sloppy
\let\oldus\_ \renewcommand{\_}{\oldus\allowbreak}
\newtheorem{theorem}{Theorem}
\newtheorem{lemma}[theorem]{Lemma}
\theoremstyle{definition}
\newtheorem{definition}[theorem]{Definition}
\title{Disproof of conjecture 00000009984:\\ squared irreducible dimensions need not divide $\dim H$}
\author{Jackmeson1}
\date{}
\begin{document}
\maketitle

\section{The conjecture}
\emph{English.} Definition: Semisimple Hopf algebras: finite-dimensional semisimple coalgebra
structures. Conjecture: The sum of squared dimensions of irreducible modules of a semisimple Hopf
algebra factors as the unique factorization of the dimension, and each squared factor divides the
dimension; the strengthened uniqueness of the factorization excludes alternating-type exceptions,
and the exception list is empty. (Hopf module dimension square factorization)

\emph{Chinese text} (as in the repository): the sum of the squares of the dimensions of the
irreducible modules decomposes as the unique decomposition of the dimension, and \emph{each square
factor divides the dimension}; the exception list is empty.

\section{Reading}
The conjecture is a conjunction. We refute one conjunct, the \emph{divisibility clause}:
\begin{quote}
(D) For every finite-dimensional semisimple Hopf algebra $A$ over $\mathbb{C}$ and every
irreducible (simple) left $A$-module $M$, the square $(\dim_{\mathbb{C}} M)^2$ divides
$\dim_{\mathbb{C}} A$.
\end{quote}
Here ``squared factor'' is read as the square $d_i^2$ of an irreducible dimension $d_i$, i.e.\ a
summand of the ``sum of squared dimensions'' $\sum_i d_i^2$. ``The exception list is empty'' means
(D) is asserted for every such $A$, with no exceptions. Since the conjunction contains (D), a
counterexample to (D) refutes the conjecture as read. The first conjunct ($\sum_i d_i^2=\dim A$)
is not used.

\emph{Not refuted:} the different reading ``each irreducible dimension $d_i$ (unsquared) divides
$\dim A$''. Our example satisfies that reading ($d_i\in\{1,2\}$, $\dim A=6$), and this submission
makes no claim about it.

\section{Definitions}
\begin{itemize}
\item A \emph{Hopf algebra} over $\mathbb{C}$ is a bialgebra $(A,m,u,\Delta,\varepsilon)$ with an
antipode $S$ satisfying $m(S\otimes\mathrm{id})\Delta=u\varepsilon=m(\mathrm{id}\otimes S)\Delta$.
In Lean this is Mathlib's class \texttt{HopfAlgebra} \texttt{\char`\\C} \texttt{A}.
\item $A$ is \emph{semisimple} if it is semisimple as an algebra, i.e.\ every left $A$-module is a
direct sum of simple modules (Lean: \texttt{IsSemisimpleRing A}).
\item An \emph{irreducible} module is a simple left $A$-module: nonzero with no submodules other
than $0$ and itself (Lean: \texttt{IsSimpleModule A M}). Its $\mathbb{C}$-structure is the
restriction of the $A$-structure (Lean: \texttt{IsScalarTower \char`\\C A M}); we also require it
to be finite-dimensional, which only weakens (D) and so strengthens the refutation.
\end{itemize}

\section{The counterexample $A=\mathbb{C}[S_3]$}
Let $S_3$ be the symmetric group on $\{0,1,2\}$ and $A=\mathbb{C}[S_3]$ its group algebra, with
the standard Hopf structure $\Delta(g)=g\otimes g$, $\varepsilon(g)=1$, $S(g)=g^{-1}$ (Mathlib's
instance on \texttt{MonoidAlgebra}). Then:
\begin{enumerate}
\item $\dim_{\mathbb{C}} A=|S_3|=6$ (Lean: \texttt{finrank\_H}).
\item $A$ is semisimple by Maschke's theorem, since $6\neq 0$ in $\mathbb{C}$ (Lean:
\texttt{H\_semisimple}, from Mathlib's Maschke instance).
\item Every $g\in S_3$ is group-like: $\Delta(g)=g\otimes g$, $\varepsilon(g)=1$, $S(g)=g^{-1}$
(Lean: \texttt{grouplike}). Hence the coalgebra $A$ is the direct sum of the one-dimensional
subcoalgebras $\mathbb{C}g$, so $A$ also has a semisimple coalgebra structure, as in the
conjecture's definition line.
\end{enumerate}
$S_3$ is a symmetric group, so $A$ is not of ``alternating type'' under any meaning of that phrase;
in any case the conjecture says the exception list is empty.

\begin{lemma}[\texttt{sq\_dvd\_six}]
If $d\in\mathbb{N}$ and $d^2\mid 6$ then $d=1$.
\end{lemma}
\begin{proof}
$d^2\le 6$ gives $d\le 2$; $d=0$ fails since $0\nmid 6$, and $d=2$ fails since $4\nmid 6$.
\end{proof}

\begin{lemma}[\texttt{comm\_on\_simple}]
Let $N\subseteq A$ be a simple left ideal with $\dim_{\mathbb{C}}N=1$. Then
$(ab-ba)n=0$ for all $a,b\in A$ and $n\in N$.
\end{lemma}
\begin{proof}
Pick $0\neq v\in N$; then $N=\mathbb{C}v$, so $av=\alpha v$ and $bv=\beta v$ for scalars
$\alpha,\beta$. Since scalars are central in $A$, $abv=\alpha\beta v=bav$, so $(ab-ba)v=0$, and
hence $(ab-ba)(cv)=c(ab-ba)v=0$.
\end{proof}

\begin{theorem}[\texttt{H\_not\_squaredDimsDivide}]
$A=\mathbb{C}[S_3]$ violates (D).
\end{theorem}
\begin{proof}
Assume (D) for $A$. Every simple left ideal $N$ of $A$ is a simple $A$-module, so
$(\dim N)^2\mid 6$ and $\dim N=1$ by the first lemma. By the second lemma, every commutator
$c=ab-ba$ annihilates every simple left ideal. Since $A$ is a semisimple module over itself, it is
the supremum of its simple submodules (Mathlib: \texttt{IsSemisimpleModule.sSup\_simples\_eq\_top}),
and the annihilator of a supremum is the intersection of the annihilators
(\texttt{Submodule.annihilator\_iSup}). So $c$ annihilates $A$, and $c=c\cdot 1=0$. Thus $A$ is
commutative. But for the transpositions $\sigma=(0\,1)$, $\tau=(1\,2)$ we have
$\sigma\tau\neq\tau\sigma$, and distinct group elements are distinct basis vectors of $A$, so
$\sigma\tau\neq\tau\sigma$ in $A$ (Lean: \texttt{H\_not\_comm}). Contradiction.
\end{proof}
(Concretely, the offending simple module is the $2$-dimensional standard representation; the Lean
proof does not need to construct it.)

\begin{theorem}[\texttt{disproof}]
It is not true that every finite-dimensional semisimple Hopf algebra $A$ over $\mathbb{C}$
satisfies (D).
\end{theorem}

\section{Lean formalization}
File \texttt{lean/Conjecture9984/Basic.lean} (Lean 4, Mathlib v4.33.1, 127 lines). Definition:
\begin{verbatim}
def SquaredDimsDivide (A : Type) [Ring A] [HopfAlgebra C A] : Prop :=
  forall (M : Type) [AddCommGroup M] [Module A M] [Module C M]
    [IsScalarTower C A M] [IsSimpleModule A M] [FiniteDimensional C M],
    (finrank C M) ^ 2 | finrank C A
\end{verbatim}
Main theorem:
\begin{verbatim}
theorem disproof :
    not (forall (A : Type) [Ring A] [HopfAlgebra C A] [IsSemisimpleRing A]
      [FiniteDimensional C A], SquaredDimsDivide A)
\end{verbatim}
(here \texttt{C} is $\mathbb{C}$). Supporting theorems: \texttt{finrank\_H}, \texttt{H\_semisimple},
\texttt{grouplike}, \texttt{H\_not\_comm}, \texttt{sq\_dvd\_six}, \texttt{comm\_on\_simple},
\texttt{H\_not\_squaredDimsDivide}.

\emph{Axiom audit.} \texttt{\#print axioms C9984.disproof} (likewise \texttt{grouplike},
\texttt{finrank\_H}) reports only \texttt{propext}, \texttt{Classical.choice},
\texttt{Quot.sound}. There is no \texttt{sorry} and no \texttt{native\_decide}.
\texttt{lake build} succeeds.

\end{document}
